\documentclass[a4paper,12pt,twoside, openright]{memoir}
\usepackage[left=4cm,right=3cm]{geometry}

% Castellano
\usepackage[spanish,es-tabla]{babel}
\selectlanguage{spanish}
\usepackage[utf8]{inputenc}
\usepackage[T1]{fontenc}
\usepackage{lmodern} % Scalable font
\usepackage{microtype}
\usepackage{placeins}

\RequirePackage{booktabs}
\RequirePackage[table]{xcolor}
\RequirePackage{xtab}
\RequirePackage{multirow}

% Links
\PassOptionsToPackage{hyphens}{url}\usepackage[colorlinks]{hyperref}
\hypersetup{
	allcolors = {red}
}

% Ecuaciones
\usepackage{amsmath}

% Rutas de fichero / paquete
\newcommand{\ruta}[1]{{\sffamily #1}}

% Párrafos
\nonzeroparskip

% Huérfanas y viudas
\widowpenalty100000
\clubpenalty100000

% Imagenes
\usepackage{graphicx}
\newcommand{\imagen}[2]{
	\begin{figure}[!h]
		\centering
		\includegraphics[width=0.9\textwidth]{#1}
		\caption{#2}\label{fig:#1}
	\end{figure}
	\FloatBarrier
}

\newcommand{\imagenflotante}[2]{
	\begin{figure}%[!h]
		\centering
		\includegraphics[width=0.9\textwidth]{#1}
		\caption{#2}\label{fig:#1}
	\end{figure}
}



% El comando \figura nos permite insertar figuras comodamente, y utilizando
% siempre el mismo formato. Los parametros son:
% 1 -> Porcentaje del ancho de página que ocupará la figura (de 0 a 1)
% 2 --> Fichero de la imagen
% 3 --> Texto a pie de imagen
% 4 --> Etiqueta (label) para referencias
% 5 --> Opciones que queramos pasarle al \includegraphics
% 6 --> Opciones de posicionamiento a pasarle a \begin{figure}
\newcommand{\figuraConPosicion}[6]{%
  \setlength{\anchoFloat}{#1\textwidth}%
  \addtolength{\anchoFloat}{-4\fboxsep}%
  \setlength{\anchoFigura}{\anchoFloat}%
  \begin{figure}[#6]
    \begin{center}%
      \Ovalbox{%
        \begin{minipage}{\anchoFloat}%
          \begin{center}%
            \includegraphics[width=\anchoFigura,#5]{#2}%
            \caption{#3}%
            \label{#4}%
          \end{center}%
        \end{minipage}
      }%
    \end{center}%
  \end{figure}%
}

%
% Comando para incluir imágenes en formato apaisado (sin marco).
\newcommand{\figuraApaisadaSinMarco}[5]{%
  \begin{figure}%
    \begin{center}%
    \includegraphics[angle=90,height=#1\textheight,#5]{#2}%
    \caption{#3}%
    \label{#4}%
    \end{center}%
  \end{figure}%
}
% Para las tablas
\newcommand{\otoprule}{\midrule [\heavyrulewidth]}
%
% Nuevo comando para tablas pequeñas (menos de una página).
\newcommand{\tablaSmall}[5]{%
 \begin{table}
  \begin{center}
   \rowcolors {2}{gray!35}{}
   \begin{tabular}{#2}
    \toprule
    #4
    \otoprule
    #5
    \bottomrule
   \end{tabular}
   \caption{#1}
   \label{tabla:#3}
  \end{center}
 \end{table}
}

%
% Nuevo comando para tablas pequeñas (menos de una página).
\newcommand{\tablaSmallSinColores}[5]{%
 \begin{table}[H]
  \begin{center}
   \begin{tabular}{#2}
    \toprule
    #4
    \otoprule
    #5
    \bottomrule
   \end{tabular}
   \caption{#1}
   \label{tabla:#3}
  \end{center}
 \end{table}
}

\newcommand{\tablaApaisadaSmall}[5]{%
\begin{landscape}
  \begin{table}
   \begin{center}
    \rowcolors {2}{gray!35}{}
    \begin{tabular}{#2}
     \toprule
     #4
     \otoprule
     #5
     \bottomrule
    \end{tabular}
    \caption{#1}
    \label{tabla:#3}
   \end{center}
  \end{table}
\end{landscape}
}

%
% Nuevo comando para tablas grandes con cabecera y filas alternas coloreadas en gris.
\newcommand{\tabla}[6]{%
  \begin{center}
    \tablefirsthead{
      \toprule
      #5
      \otoprule
    }
    \tablehead{
      \multicolumn{#3}{l}{\small\sl continúa desde la página anterior}\\
      \toprule
      #5
      \otoprule
    }
    \tabletail{
      \hline
      \multicolumn{#3}{r}{\small\sl continúa en la página siguiente}\\
    }
    \tablelasttail{
      \hline
    }
    \bottomcaption{#1}
    \rowcolors {2}{gray!35}{}
    \begin{xtabular}{#2}
      #6
      \bottomrule
    \end{xtabular}
    \label{tabla:#4}
  \end{center}
}

%
% Nuevo comando para tablas grandes con cabecera.
\newcommand{\tablaSinColores}[6]{%
  \begin{center}
    \tablefirsthead{
      \toprule
      #5
      \otoprule
    }
    \tablehead{
      \multicolumn{#3}{l}{\small\sl continúa desde la página anterior}\\
      \toprule
      #5
      \otoprule
    }
    \tabletail{
      \hline
      \multicolumn{#3}{r}{\small\sl continúa en la página siguiente}\\
    }
    \tablelasttail{
      \hline
    }
    \bottomcaption{#1}
    \begin{xtabular}{#2}
      #6
      \bottomrule
    \end{xtabular}
    \label{tabla:#4}
  \end{center}
}

%
% Nuevo comando para tablas grandes sin cabecera.
\newcommand{\tablaSinCabecera}[5]{%
  \begin{center}
    \tablefirsthead{
      \toprule
    }
    \tablehead{
      \multicolumn{#3}{l}{\small\sl continúa desde la página anterior}\\
      \hline
    }
    \tabletail{
      \hline
      \multicolumn{#3}{r}{\small\sl continúa en la página siguiente}\\
    }
    \tablelasttail{
      \hline
    }
    \bottomcaption{#1}
  \begin{xtabular}{#2}
    #5
   \bottomrule
  \end{xtabular}
  \label{tabla:#4}
  \end{center}
}



\definecolor{cgoLight}{HTML}{EEEEEE}
\definecolor{cgoExtralight}{HTML}{FFFFFF}

%
% Nuevo comando para tablas grandes sin cabecera.
\newcommand{\tablaSinCabeceraConBandas}[5]{%
  \begin{center}
    \tablefirsthead{
      \toprule
    }
    \tablehead{
      \multicolumn{#3}{l}{\small\sl continúa desde la página anterior}\\
      \hline
    }
    \tabletail{
      \hline
      \multicolumn{#3}{r}{\small\sl continúa en la página siguiente}\\
    }
    \tablelasttail{
      \hline
    }
    \bottomcaption{#1}
    \rowcolors[]{1}{cgoExtralight}{cgoLight}

  \begin{xtabular}{#2}
    #5
   \bottomrule
  \end{xtabular}
  \label{tabla:#4}
  \end{center}
}



\graphicspath{ {./img/} }

% Capítulos
\chapterstyle{bianchi}
\newcommand{\capitulo}[2]{
	\setcounter{chapter}{#1}
	\setcounter{section}{0}
	\setcounter{figure}{0}
	\setcounter{table}{0}
	\chapter*{#2}
	\addcontentsline{toc}{chapter}{#2}
	\markboth{#2}{#2}
}

% Apéndices
\renewcommand{\appendixname}{Apéndice}
\renewcommand*\cftappendixname{\appendixname}

\newcommand{\apendice}[1]{
	%\renewcommand{\thechapter}{A}
	\chapter{#1}
}

\renewcommand*\cftappendixname{\appendixname\ }

% Formato de portada
\makeatletter
\usepackage{xcolor}
\newcommand{\tutor}[1]{\def\@tutor{#1}}
\newcommand{\course}[1]{\def\@course{#1}}
\definecolor{cpardoBox}{HTML}{E6E6FF}
\def\maketitle{
  \null
  \thispagestyle{empty}
  % Cabecera ----------------
\noindent\includegraphics[width=\textwidth]{cabecera}\vspace{1cm}%
  \vfill
  % Título proyecto y escudo informática ----------------
  \colorbox{cpardoBox}{%
    \begin{minipage}{.8\textwidth}
      \vspace{.5cm}\Large
      \begin{center}
      \textbf{TFG del Grado en Ingeniería Informática}\vspace{.6cm}\\
      \textbf{\LARGE\@title{}}
      \end{center}
      \vspace{.2cm}
    \end{minipage}

  }%
  \hfill\begin{minipage}{.20\textwidth}
    \includegraphics[width=\textwidth]{escudoInfor}
  \end{minipage}
  \vfill
  % Datos de alumno, curso y tutores ------------------
  \begin{center}%
  {%
    \noindent\LARGE
    Presentado por \@author{}\\ 
    en Universidad de Burgos --- \@date{}\\
    Tutores: \@tutor{}\\
  }%
  \end{center}%
  \null
  \cleardoublepage
  }
\makeatother

\newcommand{\nombre}{José Daniel Ballester Delgado} %%% cambio de comando

% Datos de portada
\title{Modern Data Stack para analizar la situación epidemiológica}
\author{\nombre}
\tutor{José Ignacio Santos Martín y José Manuel Galán Ordax}
\date{\today}

\begin{document}

\maketitle


\newpage\null\thispagestyle{empty}\newpage


%%%%%%%%%%%%%%%%%%%%%%%%%%%%%%%%%%%%%%%%%%%%%%%%%%%%%%%%%%%%%%%%%%%%%%%%%%%%%%%%%%%%%%%%
\thispagestyle{empty}


\noindent\includegraphics[width=\textwidth]{cabecera}\vspace{1cm}

\noindent D. José Ignacio Santos Martín y D. José Manuel Galán Ordax, profesores del departamento de Ingeniería Civil, área de Ingeniería de Organización.

\noindent Exponen:

\noindent Que el alumno D. \nombre, con DNI 71565431B, ha realizado el Trabajo final de Grado en Ingeniería Informática titulado Modern Data Stack para analizar la situación epidemiológica. 

\noindent Y que dicho trabajo ha sido realizado por el alumno bajo la dirección del que suscribe, en virtud de lo cual se autoriza su presentación y defensa.

\begin{center} %\large
En Burgos, {\large \today}
\end{center}

\vfill\vfill\vfill

% Author and supervisor
\begin{minipage}{0.45\textwidth}
\begin{flushleft} %\large
Vº. Bº. del Tutor:\\[2cm]
D. José Ignacio Santos Martín
\end{flushleft}
\end{minipage}
\hfill
\begin{minipage}{0.45\textwidth}
\begin{flushleft} %\large
Vº. Bº. del co-tutor:\\[2cm]
D. José Manuel Galán Ordax
\end{flushleft}
\end{minipage}
\hfill

\vfill

% para casos con solo un tutor comentar lo anterior
% y descomentar lo siguiente
%Vº. Bº. del Tutor:\\[2cm]
%D. nombre tutor


\newpage\null\thispagestyle{empty}\newpage




\frontmatter

% Abstract en castellano
\renewcommand*\abstractname{Resumen}
\begin{abstract}

La gran diversidad de fuentes de información en internet provoca desconfianza a la hora de consultar datos, ya que en muchas ocasiones es falsa o no está verificada. Por tanto este trabajo recopilará información actualizada sobre la situación de la pandemia de las principales fuentes oficiales, y brindará una opción visual e interactiva que proporciona el acceso a esta información veraz, sobre el estado de la pandemia causada por el Covid-19 en España.

Y todo ello se desarrollará utilizando un nuevo paradigma a la hora de construir la arquitectura que da soporte al ciclo de vida de los datos, el Modern Data Stack, que haciendo uso de la tecnología en la nube, trae muchas ventajas con respecto a la arquitectura tradicional.
\end{abstract}

\renewcommand*\abstractname{Descriptores}
\begin{abstract}
Análisis de datos, Covid-19, Business Intelligence, Data Warehouse, KPI, Nube, OLAP, ELT, Modern Data Stack, Analytics engineer.

\end{abstract}

\clearpage

% Abstract en inglés
\renewcommand*\abstractname{Abstract}
\begin{abstract}
The great diversity of information sources on the internet causes mistrust when consulting data, since it is often false or unverified. Therefore, this work will compile updated information of the pandemic's situation, from the main official sources and will provide a visual and interactive option that will give access to this truthful information, about the state of the pandemic caused by Covid-19 in Spain.

And all this will be developed using a new architecture paradigm that supports the data life cycle, the Modern Data Stack, which, by making use of cloud technology, brings many advantages in comparison to the traditional architecture.
\end{abstract}

\renewcommand*\abstractname{Keywords}
\begin{abstract}
Data analysis, Covid-19, Business Intelligence, Data Warehouse, KPI, Cloud, OLAP, ELT, Modern Data Stack, Analytics engineer
\end{abstract}

\clearpage

% Indices
\tableofcontents

\clearpage

\listoffigures

\clearpage

\listoftables
\clearpage

\mainmatter
\capitulo{1}{Introducción}
La gran diversidad de fuentes de información en internet provoca desconfianza a la hora de consultar datos, ya que en muchas ocasiones la información es falsa o no está verificada.

Este trabajo tiene como objetivo brindar una solución visual e interactiva que facilite el acceso a información de la situación del Covid-19 en España.

Como respuesta a este problema se prevé desarrollar un cuadro de mandos que recopile los principales datos de las fuentes oficiales encargadas de publicar información sobre la situación actual del Covid-19 en el país, de forma que el usuario tenga una opción centralizada y fiable para consultar la información que desea y poder observar insights sobre ésta. 

Asimismo, otra novedad que aporta con respecto a las alternativas más clásicas, es la interacción del usuario con los datos, permitiendo una modificación dinámica de la visualización mediante el uso de filtros, de manera que permita obtener información específica y más detallada sobre un dato o un conjunto de datos.

Estos análisis son complicados con los sistemas transaccionales tradicionales (OLTP), porque la información está dispersa en muchas tablas y no están dirigidas al análisis. 

Debido a estos problemas, se ha introducido el uso de Business Intelligence y los sistemas analíticos (OLAP), donde la información se integra en un repositorio optimizado, para responder preguntas estratégicas. 

El Business Intelligence (BI) se conoce como la capacidad de transformar los datos en información y la información en conocimiento. Se logra a través del almacenamiento y procesamiento de grandes cantidades de datos para que la información pueda ser analizada.

Otro aspecto importante, es que en este TFG se pretende utilizar el Modern Data Stack, un nuevo paradigma a la hora de montar la arquitectura que da soporte al ciclo de vida de los datos, gracias a las ventajas que tiene montar esta estructura en la nube.

El término “data stack” se origina en “technology stack”, la combinación de diferentes tecnologías, que los ingenieros de software combinan para crear productos y servicios. 

Los data stack están diseñados específicamente para respaldar el almacenamiento, la administración y el acceso a los datos y se usan para aprovechar los datos en la toma de decisiones estratégicas. 

El Modern Data Stack (MDS) consiste en un conjunto flexible de tecnologías que ayudan a almacenar, administrar y aprender de los datos.

La nube ha transformado la industria del software y la forma en que las empresas construyen sus productos. Hoy, podemos configurar un technology stack en una fracción del tiempo menor y de forma más económica. Y no sorprende que estas transformaciones hayan allanado el camino para el Modern Data Stack (MDS) y un cambio de paradigma en este tipo de soluciones. 

\section{Estructura de la memoria}\label{estructura-de-la-memoria}
La memoria está formada por la siguiente estructura:
\begin{itemize}
\tightlist
\item
  \textbf{Introducción:} concisa presentación del problema a solucionar junto con la resolución sugerida. Se incluye la estructura de la memoria y los materiales adjuntos.
\item
  \textbf{Objetivos del proyecto:} definición de los objetivos que
  pretende alcanzar el proyecto.
\item
  \textbf{Conceptos teóricos:} breve desarrollo de los conceptos
  teóricos fundamentales para la correcta interpretación del resultado planteado.
\item
  \textbf{Técnicas y herramientas:} exposición de métodos y
  herramientas empleadas para organización y desarrollo del proyecto.
\item
  \textbf{Aspectos relevantes del desarrollo:} muestra de fases
  relevantes a destacar durante el transcurso de la elaboración del proyecto.
\item
  \textbf{Trabajos relacionados:} estado del arte en el campo de las arquitecturas de datos para el análisis de datos mediante reporting.
\item
  \textbf{Conclusiones y líneas de trabajo futuras:} conclusiones
  extraídas de los resultados obtenidos a partir de la realización del proyecto y opciones de mejora o
  ampliación de la resolución planteada.
\end{itemize}

\section{Materiales adjuntos}\label{materiales-adjuntos}
\begin{itemize}	
	\item \textbf{Anexos:}
	\begin{itemize}
	\tightlist
	\item
  		\textbf{Plan del proyecto software:} organización y potencialidad del proyecto.
	\item
  		\textbf{Especificación de requisitos del software:} se incluye la
  		fase de análisis; los objetivos generales, el catálogo de requisitos
  		del sistema y la especificación de requisitos funcionales y no
  		funcionales.
	\item
  		\textbf{Especificación de diseño:} se expone la fase de diseño; el
  		ámbito del software, el diseño de datos, el diseño procedimental y el
  		diseño arquitectónico.
	\item
  		\textbf{Manual del programador:} comprende los aspectos más relevantes
  		relacionados con el código fuente (estructura, compilación,
  		instalación, ejecución, pruebas, etc.).
	\item
  		\textbf{Manual de usuario:} guía de usuario para el buen uso y empleo adecuado de la aplicación.
	\end{itemize}
	
	\item \textbf{Repositorio Github:} contiene la memoria y los anexos, además del archivo del cuadro de mandos, los accesos los sercivios y herramientas cloud utlilizados y los scripts de dbt, Snowflake y Phyton del proyecto.
	
\end{itemize}

\capitulo{2}{Objetivos del proyecto}

Este apartado explica de forma precisa y concisa cuales son los objetivos que se persiguen con la realización del proyecto. Se puede distinguir entre los objetivos marcados por los requisitos del software a construir y los objetivos de carácter técnico que plantea a la hora de llevar a la práctica el proyecto.

\capitulo{3}{Conceptos teóricos}

En aquellos proyectos que necesiten para su comprensión y desarrollo de unos conceptos teóricos de una determinada materia o de un determinado dominio de conocimiento, debe existir un apartado que sintetice dichos conceptos.

Algunos conceptos teóricos de \LaTeX \footnote{Créditos a los proyectos de Álvaro López Cantero: Configurador de Presupuestos y Roberto Izquierdo Amo: PLQuiz}.

\section{Secciones}

Las secciones se incluyen con el comando section.

\subsection{Subsecciones}

Además de secciones tenemos subsecciones.

\subsubsection{Subsubsecciones}

Y subsecciones. 


\section{Referencias}

Las referencias se incluyen en el texto usando cite \cite{wiki:latex}. Para citar webs, artículos o libros \cite{koza92}.


\section{Imágenes}

Se pueden incluir imágenes con los comandos standard de \LaTeX, pero esta plantilla dispone de comandos propios como por ejemplo el siguiente:

\imagen{escudoInfor}{Autómata para una expresión vacía}



\section{Listas de items}

Existen tres posibilidades:

\begin{itemize}
	\item primer item.
	\item segundo item.
\end{itemize}

\begin{enumerate}
	\item primer item.
	\item segundo item.
\end{enumerate}

\begin{description}
	\item[Primer item] más información sobre el primer item.
	\item[Segundo item] más información sobre el segundo item.
\end{description}
	
\begin{itemize}
\item 
\end{itemize}

\section{Tablas}

Igualmente se pueden usar los comandos específicos de \LaTeX o bien usar alguno de los comandos de la plantilla.

\tablaSmall{Herramientas y tecnologías utilizadas en cada parte del proyecto}{l c c c c}{herramientasportipodeuso}
{ \multicolumn{1}{l}{Herramientas} & App AngularJS & API REST & BD & Memoria \\}{ 
HTML5 & X & & &\\
CSS3 & X & & &\\
BOOTSTRAP & X & & &\\
JavaScript & X & & &\\
AngularJS & X & & &\\
Bower & X & & &\\
PHP & & X & &\\
Karma + Jasmine & X & & &\\
Slim framework & & X & &\\
Idiorm & & X & &\\
Composer & & X & &\\
JSON & X & X & &\\
PhpStorm & X & X & &\\
MySQL & & & X &\\
PhpMyAdmin & & & X &\\
Git + BitBucket & X & X & X & X\\
Mik\TeX{} & & & & X\\
\TeX{}Maker & & & & X\\
Astah & & & & X\\
Balsamiq Mockups & X & & &\\
VersionOne & X & X & X & X\\
} 

\capitulo{4}{Técnicas y herramientas}

Esta parte de la memoria tiene como objetivo presentar las técnicas metodológicas y las herramientas de desarrollo que se han utilizado para llevar a cabo el proyecto. Si se han estudiado diferentes alternativas de metodologías, herramientas, bibliotecas se puede hacer un resumen de los aspectos más destacados de cada alternativa, incluyendo comparativas entre las distintas opciones y una justificación de las elecciones realizadas. 
No se pretende que este apartado se convierta en un capítulo de un libro dedicado a cada una de las alternativas, sino comentar los aspectos más destacados de cada opción, con un repaso somero a los fundamentos esenciales y referencias bibliográficas para que el lector pueda ampliar su conocimiento sobre el tema.



\capitulo{5}{Aspectos relevantes del desarrollo del proyecto}

Este apartado pretende recoger los aspectos más interesantes del desarrollo del proyecto, comentados por los autores del mismo.
Debe incluir desde la exposición del ciclo de vida utilizado, hasta los detalles de mayor relevancia de las fases de análisis, diseño e implementación.
Se busca que no sea una mera operación de copiar y pegar diagramas y extractos del código fuente, sino que realmente se justifiquen los caminos de solución que se han tomado, especialmente aquellos que no sean triviales.
Puede ser el lugar más adecuado para documentar los aspectos más interesantes del diseño y de la implementación, con un mayor hincapié en aspectos tales como el tipo de arquitectura elegido, los índices de las tablas de la base de datos, normalización y desnormalización, distribución en ficheros3, reglas de negocio dentro de las bases de datos (EDVHV GH GDWRV DFWLYDV), aspectos de desarrollo relacionados con el WWW...
Este apartado, debe convertirse en el resumen de la experiencia práctica del proyecto, y por sí mismo justifica que la memoria se convierta en un documento útil, fuente de referencia para los autores, los tutores y futuros alumnos.

\capitulo{6}{Trabajos relacionados}

Este apartado es un estado del arte. Incluye un pequeño resumen comentado de los trabajos y proyectos ya realizados en el campo del proyecto en curso. 

Al empezar con el TFG busqué trabajos relacionados con el mío para ver cómo estaba el estado del arte, pero como la mayoría de los trabajos de este tipo los hacen empresas con datos sensibles no pude encontrar apenas. Lo que sí encontré fueron TFGs de otras universidades públicos, los cuales me sirvieron para poder compararlos con lo que yo tenía intención de hacer.

\section{Análisis y visualización de datos abiertos de 
carácter informativo para el alumnado de la 
Universidad de Valladolid}

Se trata de un trabajo de fin de grado realizado por Nadia Pérez Faúndez en el año 2021 \cite{TFG1}.

Su proyecto se centra en la disposición de datos abiertos con el fin de informar a los estudiantes universitarios a través de una aplicación. 

Aunque el proceso de desarrollo no es el mismo que el planteado en este proyecto, ya que usa la arquitectura tradicional con ETL, se puede ver como el objetivo final se asemeja, ya que también pretende la extracción de datos fiables de fuentes oficiales, en este caso de distintas universidades públicas de Castilla y León, con la finalidad de proporcionar la visualización de la información obtenida facilitando al usuario el acceso e interacción con los datos.


\section{Aplicación de herramientas de Business Intelligence en datos del entorno de Salud}
Se trata de un trabajo de fin de grado realizado por Uxue Ayechu Abendaño en el año 2017 \cite{TFG2}.

Su proyecto se centra en la explotación de la información sanitaria de Navarra a través del Business Intelligence con el objetivo de ayudar a los sanitarios a partir de los datos obtendidos, con su visualización a través de la herramienta Tableau.

En su caso, se usa la arquitectura tradicional con ETL, a diferencia de este proyecto, en cuyo proceso de desarrollo usa el Modern Data Stack, estando todo almacenado en la nube y usando los procesos ELT.

\section{Business Intelligence y el análisis predictivo: COVID 19}
Se trata de un trabajo de fin de grado realizado por Javier López Navas en el año 2020 \cite{TFG3}.

Su proyecto se centra en elaborar un estudio de Business Intelligence del Covid-19 a través del uso de la herramienta Power Bi de Microsoft, mediante la interpretación de datos. 

En este caso se emplea la misma herramienta para la visualización de la información, aunque la principal diferencia es en el proceso de desarrollo, ya que en este proyecto se usa la arquitectura tradicional con ETL y en el proceso de desarrollo de este proyecto se usa el Modern Data Stack, estando todo almacenado en la nube y usando los procesos ELT.

\section{Mejoras de mi proyecto con respecto a los analizados en este apartado}
Gracias a utilizar este nuevo paradigma de arquitectura de datos, mi proyecto tiene unas ventajas que no tienen proyectos anteriores:
\begin{itemize}
    \item Mayor velocidad de procesamiento en los datos.
    \item No necesitar instalaciones al usar servicios en la nube.
    \item Implementar practicas de Ingeniería del sofware gracias a dbt como: modularidad, portabilidad, integación, distribución continua, documentación, control de versiones y testeo.
    \item Actualización diaria automática de los datos en todas las fases de nuestra arquitectura, con su correspondiente testeo.
    \item Gracias al enfoque ELT en lugar de ETL obtenemos menos latencia, mayor escalabilidad y menor procesamiento.
\end{itemize}
\capitulo{7}{Conclusiones y Líneas de trabajo futuras}

Conclusiones derivadas del desarrollo del proyecto.  

\section{Conclusiones}
Una vez finalizado el proyecto, se pueden considerar una serie de conclusiones basadas en su evolución:
\begin{itemize}
	\item Se ha podido comprobar que los objetivos generales se han cumplido de manera satisfactoria, ya que finalmente se ha logrado el desarrollo de una arquitectura de datos que recopila la información de las principales fuentes oficiales sobre el estado de la pandemia en España, mediante su extracción, carga y transformación en la nube. Y, además, de disponer toda la información recopilada almacenada en la nube, se ha facilitado al usuario la interacción con los datos.
	
	\item El proyecto ha requerido el uso de nuevos conocimientos sobre distintos conceptos y diversas herramientas.
	Respecto a los nuevos conceptos aprendidos puedo considerar que, en los que más he tenido que profundizar, son aquellos conceptos relacionados con el nuevo paradigma \emph{modern data stack}, entre los que se incluyen ELT y \emph{analytics engineer}. El aprendizaje de todos estos conceptos ha resultado muy útil debido al aumento de uso de esta nueva arquitectura que están teniendo las empresas españolas, gracias a las ventajas que ofrece.

	\item La evolución del proyecto ha sido controlada a través del desarrollo ágil de Scrum a partir de GitHub, lo que ha facilitado la organización y aceleración de la entrega final del proyecto. Gracias a la elaboración del proyecto en sprints se ha podido observar el incremento de valor y, a su vez corregir los errores localizados en sprints anteriores.
\end{itemize}

 \section{Líneas de trabajo futuras}
 Algunas de las líneas de trabajo futuras que se podrían considerar son:
\begin{itemize}
    \item Conseguir una automatización completa en la extracción de datos.
    \item Implementar algún análisis predictivo, mediante técnicas estadísticas de modelización, aprendizaje automático y minería de datos.
    \item Conseguir una versión premium de PowerBI para poder hacer público la URL del cuadro de mandos.
    \item Añadir más gráficos para la visualización de datos con nuevos KPIs en la herramienta empleada {\emph{Microsoft Power BI}}.
\end{itemize}
 





\bibliographystyle{plain}
\bibliography{bibliografia}

\end{document}
